\documentclass[10pt,a4paper]{amsart}
\usepackage{amsmath,amssymb,amsthm,graphicx}
% For commutative diagrams we use Xy-pic
\usepackage[all]{xy}

% We use 2cell for 2-commutative diagrams.
\xyoption{2cell}
% Theorem environments.
%
\theoremstyle{plain}
\newtheorem{theorem}[subsection]{Theorem}
\newtheorem{proposition}[subsection]{Proposition}
\newtheorem{lemma}[subsection]{Lemma}

\theoremstyle{definition}
\newtheorem{definition}[subsection]{Definition}
\newtheorem{example}[subsection]{Example}
\newtheorem{exercise}[subsection]{Exercise}
\newtheorem{situation}[subsection]{Situation}

\theoremstyle{remark}
\newtheorem{remark}[subsection]{Remark}
\newtheorem{remarks}[subsection]{Remarks}

\numberwithin{equation}{subsection}

% Macros
%
\def\lim{\mathop{\mathrm{lim}}\nolimits}
\def\colim{\mathop{\mathrm{colim}}\nolimits}
\def\Spec{\mathop{\mathrm{Spec}}}
\def\Hom{\mathop{\mathrm{Hom}}\nolimits}
\def\Ext{\mathop{\mathrm{Ext}}\nolimits}
\def\SheafHom{\mathop{\mathcal{H}\!\mathit{om}}\nolimits}
\def\SheafExt{\mathop{\mathcal{E}\!\mathit{xt}}\nolimits}
\def\Sch{\mathit{Sch}}
\def\Mor{\mathop{\mathrm{Mor}}\nolimits}
\def\Ob{\mathop{\mathrm{Ob}}\nolimits}
\def\Sh{\mathop{\mathit{Sh}}\nolimits}
\def\NL{\mathop{N\!L}\nolimits}
\def\CH{\mathop{\mathrm{CH}}\nolimits}
\def\proetale{{pro\text{-}\acute{e}tale}}
\def\etale{{\acute{e}tale}}
\def\QCoh{\mathit{QCoh}}
\def\Ker{\mathop{\mathrm{Ker}}}
\def\Im{\mathop{\mathrm{Im}}}
\def\Coker{\mathop{\mathrm{Coker}}}
\def\Coim{\mathop{\mathrm{Coim}}}

% Boxtimes
%
\DeclareMathSymbol{\boxtimes}{\mathbin}{AMSa}{"02}

%
% Macros for moduli stacks/spaces
%
\def\QCohstack{\mathcal{QC}\!\mathit{oh}}
\def\Cohstack{\mathcal{C}\!\mathit{oh}}
\def\Spacesstack{\mathcal{S}\!\mathit{paces}}
\def\Quotfunctor{\mathrm{Quot}}
\def\Hilbfunctor{\mathrm{Hilb}}
\def\Curvesstack{\mathcal{C}\!\mathit{urves}}
\def\Polarizedstack{\mathcal{P}\!\mathit{olarized}}
\def\Complexesstack{\mathcal{C}\!\mathit{omplexes}}
% \Pic is the operator that assigns to X its picard group, usage \Pic(X)
% \Picardstack_{X/B} denotes the Picard stack of X over B
% \Picardfunctor_{X/B} denotes the Picard functor of X over B
\def\Pic{\mathop{\mathrm{Pic}}\nolimits}
\def\Picardstack{\mathcal{P}\!\mathit{ic}}
\def\Picardfunctor{\mathrm{Pic}}
\def\Deformationcategory{\mathcal{D}\!\mathit{ef}}
\newenvironment{reference}{}{}
\newenvironment{slogan}{}{}
\newenvironment{history}{}{}
\usepackage{cleveref}
\crefname{originallemma}{original lemma}{original lemmas}
\Crefname{originallemma}{Original lemma}{Original lemmas}
\makeatletter
\newlabel{lemma-lift-complex}{{0DIT}{source}{}{sourcepointer-1}{}}
\newlabel{lemma-lift-complex@cref}{{[originallemma][0][]0DIT}{[1][0][]source}{}{}{}}
\makeatother
\crefname{originallemma}{original lemma}{original lemmas}
\Crefname{originallemma}{Original lemma}{Original lemmas}
\makeatletter
\newlabel{derived-lemma-negative-vanishing}{{05TC}{source}{}{sourcepointer-2}{}}
\newlabel{derived-lemma-negative-vanishing@cref}{{[originallemma][0][]05TC}{[1][0][]source}{}{}{}}
\makeatother
\crefname{originallemma}{original lemma}{original lemmas}
\Crefname{originallemma}{Original lemma}{Original lemmas}
\makeatletter
\newlabel{homology-lemma-four-lemma}{{05QA}{source}{}{sourcepointer-3}{}}
\newlabel{homology-lemma-four-lemma@cref}{{[originallemma][0][]05QA}{[1][0][]source}{}{}{}}
\makeatother
\title{Vanishing of the canonical Ext-two class characterizes lifting a bounded-above complex across a square-zero thickening}
\author{The Stacks Project Authors}
\date{stacks-project:defos.tex:local-snapshot-20261011; GNU FDL 1.2 or later}
\begin{document}
\maketitle
\noindent\textit{Editorial note.} The task retains the original resolution-based tensor conventions from the canonical-class construction. It does not silently assumeK\_0 is a single flat module or perfect.\par\medskip
\noindent\textit{Allowed standard prerequisites (editorial).}\begin{itemize}
\item Explicit prerequisite boundary. Sites and sheaves of rings, bounded-above derived categories, flat resolutions, localization of the homotopy category, mapping cones and the four lemma. The canonical class and induced-map support are retained rather than assumed to give the lifting conclusion.
\end{itemize}\medskip
\ifdefined\hypertarget\hypertarget{sourcepointer-1}{}\fi
\noindent\textit{Source reference.} Original lemma 0DIT: focused-seeds/AGP/sources/stacks-snapshot-20261011/tags.txt, source lines 16953--16953. Stable Stacks tag for an expressly allowed foundational or notational input; the counted construction remains the two concrete core obligations. The retained local core is not replaced by this pointer. See standard\_prerequisites.\par\medskip
\ifdefined\hypertarget\hypertarget{sourcepointer-2}{}\fi
\noindent\textit{Source reference.} Original lemma 05TC: focused-seeds/AGP/sources/stacks-snapshot-20261011/tags.txt, source lines 7476--7476. Stable Stacks tag for an expressly allowed foundational or notational input; the counted construction remains the two concrete core obligations. The retained local core is not replaced by this pointer. See standard\_prerequisites.\par\medskip
\ifdefined\hypertarget\hypertarget{sourcepointer-3}{}\fi
\noindent\textit{Source reference.} Original lemma 05QA: focused-seeds/AGP/sources/stacks-snapshot-20261011/tags.txt, source lines 7369--7369. Stable Stacks tag for an expressly allowed foundational or notational input; the counted construction remains the two concrete core obligations. The retained local core is not replaced by this pointer. See standard\_prerequisites.\par\medskip
\begin{lemma}
\label{lemma-canonical-class-obstruction}
Let $\mathcal{C}$ be a site. Let $\mathcal{O} \to \mathcal{O}_0$
be a surjection of sheaves of rings whose kernel is an ideal sheaf
$\mathcal{I}$ of square zero. For every object
$K_0$ in $D^-(\mathcal{O}_0)$ the following are equivalent
\begin{enumerate}
\item  the class
$\omega(K_0) \in
\Ext^2_{\mathcal{O}_0}(K_0, K_0 \otimes_{\mathcal{O}_0} \mathcal{I})$
constructed in Lemma \ref{lemma-canonical-class} is zero,
\item there exists $K \in D^-(\mathcal{O})$ with
$K \otimes_\mathcal{O}^\mathbf{L} \mathcal{O}_0 = K_0$
in $D(\mathcal{O}_0)$.
\end{enumerate}
\end{lemma}

\begin{lemma}
\label{lemma-canonical-class}
Let $\mathcal{C}$ be a site. Let $\mathcal{O} \to \mathcal{O}_0$
be a surjection of sheaves of rings whose kernel is an ideal sheaf
$\mathcal{I}$ of square zero. For every object
$K_0$ in $D^-(\mathcal{O}_0)$ there is a canonical map
$$
\omega(K_0) :
K_0 \longrightarrow
K_0 \otimes_{\mathcal{O}_0}^\mathbf{L} \mathcal{I}[2]
$$
in $D(\mathcal{O}_0)$ such that for any map
$K_0 \to L_0$ in $D^-(\mathcal{O}_0)$ the diagram
$$
\xymatrix{
K_0 \ar[d] \ar[rr]_-{\omega(K_0)} & &
(K_0 \otimes^\mathbf{L}_{\mathcal{O}_0} \mathcal{I})[2] \ar[d] \\
L_0 \ar[rr]^-{\omega(L_0)} & &
(L_0 \otimes^\mathbf{L}_{\mathcal{O}_0} \mathcal{I})[2]
}
$$
commutes.
\end{lemma}

\begin{proof}
Represent $K_0$ by any complex
$\mathcal{K}_0^\bullet$ of $\mathcal{O}_0$-modules.
Apply Lemma \ref{lemma-lift-complex}
with $\mathcal{G}^n = 0$ for all $n$.
Denote $d : \mathcal{F}^n \to \mathcal{F}^{n + 1}$
the maps produced by the lemma. Then we see that
$d \circ d : \mathcal{F}^n \to \mathcal{F}^{n + 2}$
is zero modulo $\mathcal{I}$. Since $\mathcal{F}^n$ is flat,
we see that
$\mathcal{I}\mathcal{F}^n =
\mathcal{F}^n \otimes_{\mathcal{O}} \mathcal{I} =
\mathcal{F}^n_0 \otimes_{\mathcal{O}_0} \mathcal{I}$.
Hence we obtain a canonical map of complexes
$$
d \circ d : \mathcal{F}_0^\bullet
\longrightarrow
(\mathcal{F}_0^\bullet \otimes_{\mathcal{O}_0} \mathcal{I})[2]
$$
Since $\mathcal{F}_0^\bullet$ is a bounded above complex
of flat $\mathcal{O}_0$-modules, it is K-flat and may be used
to compute derived tensor product. Moreover, the map of
complexes $\mathcal{F}_0^\bullet \to \mathcal{K}_0^\bullet$
is a quasi-isomorphism by construction. Therefore the source and
target of the map just constructed represent $K_0$ and
$K_0 \otimes_{\mathcal{O}_0}^\mathbf{L} \mathcal{I}[2]$
and we obtain our map $\omega(K_0)$.

\medskip\noindent
Let us show that this procedure is compatible with maps of complexes.
Namely, let $\mathcal{L}_0^\bullet$ represent another object of
$D^-(\mathcal{O}_0)$ and suppose that
$$
\mathcal{K}_0^\bullet \longrightarrow \mathcal{L}_0^\bullet
$$
is a map of complexes. Apply Lemma \ref{lemma-lift-complex}
for the complex $\mathcal{L}_0^\bullet$, the flat modules
$\mathcal{F}^n$, the maps $\mathcal{F}^n \to \mathcal{F}^{n + 1}$,
and the compositions
$\mathcal{F}^n \to \mathcal{K}_0^n \to \mathcal{L}_0^n$
(we apologize for the reversal of letters used).
We obtain flat modules $\mathcal{G}^n$, maps
$\mathcal{F}^n \to \mathcal{G}^n$, maps
$\mathcal{G}^n \to \mathcal{G}^{n + 1}$, and maps
$\mathcal{G}^n \to \mathcal{L}_0^n$ with all properties
as in the lemma. Then it is clear that
$$
\xymatrix{
\mathcal{F}_0^\bullet \ar[d] \ar[r] &
(\mathcal{F}_0^\bullet \otimes_{\mathcal{O}_0} \mathcal{I})[2] \ar[d] \\
\mathcal{G}_0^\bullet \ar[r] &
(\mathcal{G}_0^\bullet \otimes_{\mathcal{O}_0} \mathcal{I})[2]
}
$$
is a commutative diagram of complexes.

\medskip\noindent
To see that $\omega(K_0)$ is well defined, suppose that we have two complexes
$\mathcal{K}_0^\bullet$ and $(\mathcal{K}'_0)^\bullet$
of $\mathcal{O}_0$-modules representing $K_0$ and two systems
$(\mathcal{F}^n, d : \mathcal{F}^n \to \mathcal{F}^{n + 1},
\mathcal{F}^n \to \mathcal{K}_0^n)$
and
$((\mathcal{F}')^n, d : (\mathcal{F}')^n \to (\mathcal{F}')^{n + 1},
(\mathcal{F}')^n \to \mathcal{K}_0^n)$
as above. Then we can choose a complex $(\mathcal{K}''_0)^\bullet$
and quasi-isomorphisms
$\mathcal{K}_0^\bullet  \to (\mathcal{K}''_0)^\bullet$
and
$(\mathcal{K}'_0)^\bullet  \to (\mathcal{K}''_0)^\bullet$
realizing the fact that both complexes represent $K_0$ in the
derived category. Next, we apply the result of the previous paragraph
to
$$
(\mathcal{K}_0)^\bullet \oplus (\mathcal{K}'_0)^\bullet
\longrightarrow
(\mathcal{K}''_0)^\bullet
$$
This produces a commutative diagram
$$
\xymatrix{
\mathcal{F}_0^\bullet \oplus (\mathcal{F}'_0)^\bullet
\ar[d] \ar[r] &
(\mathcal{F}_0^\bullet \otimes_{\mathcal{O}_0} \mathcal{I})[2] \oplus
((\mathcal{F}'_0)^\bullet \otimes_{\mathcal{O}_0} \mathcal{I})[2] \ar[d] \\
\mathcal{G}_0^\bullet \ar[r] &
(\mathcal{G}_0^\bullet \otimes_{\mathcal{O}_0} \mathcal{I})[2]
}
$$
Since the vertical arrows give quasi-isomorphisms on the summands
we conclude the desired commutativity in $D(\mathcal{O}_0)$.

\medskip\noindent
Having established well-definedness, the statement on compatibility
with maps is a consequence of the result in the second
paragraph.
\end{proof}

\begin{lemma}
\label{lemma-induced-map}
Let $(\mathcal{C}, \mathcal{O})$ be a ringed site.
Let $\alpha : K \to L$ be a map of $D^-(\mathcal{O})$.
Let $\mathcal{F}$ be a sheaf of $\mathcal{O}$-modules.
Let $n \in \mathbf{Z}$.
\begin{enumerate}
\item If $H^i(\alpha)$ is an isomorphism for $i \geq n$,
then $H^i(\alpha \otimes_\mathcal{O}^\mathbf{L} \text{id}_\mathcal{F})$
is an isomorphism for $i \geq n$.
\item If $H^i(\alpha)$ is an isomorphism for $i > n$ 
and surjective for $i = n$,
then $H^i(\alpha \otimes_\mathcal{O}^\mathbf{L} \text{id}_\mathcal{F})$
is an isomorphism for $i > n$ and surjective for $i = n$.
\end{enumerate}
\end{lemma}

\begin{proof}
Choose a distinguished triangle
$$
K \to L \to C \to K[1]
$$
In case (2) we see that $H^i(C) = 0$ for $i \geq n$.
Hence $H^i(C \otimes_\mathcal{O}^\mathbf{L} \mathcal{F}) = 0$
for $i \geq n$ by (the dual of)
Derived Categories, Lemma \ref{derived-lemma-negative-vanishing}.
This in turn shows that
$H^i(\alpha \otimes_\mathcal{O}^\mathbf{L} \text{id}_\mathcal{F})$
is an isomorphism for $i > n$ and surjective for $i = n$.
In case (1) we moreover see that $H^{n - 1}(L) \to H^{n - 1}(C)$
is surjective. Considering the diagram
$$
\xymatrix{
H^{n - 1}(L) \otimes_\mathcal{O} \mathcal{F} \ar[r] \ar[d] &
H^{n - 1}(C) \otimes_\mathcal{O} \mathcal{F} \ar@{=}[d] \\
H^{n - 1}(L \otimes_\mathcal{O}^\mathbf{L} \mathcal{F}) \ar[r] &
H^{n - 1}(C \otimes_\mathcal{O}^\mathbf{L} \mathcal{F})
}
$$
we conclude the lower horizontal arrow is surjective. Combined with what
was said before this implies that
$H^n(\alpha \otimes_\mathcal{O}^\mathbf{L} \text{id}_\mathcal{F})$
is an isomorphism.
\end{proof}

\begin{proof}
Let $K$ be as in (2). Then we can represent $K$ by a bounded above
complex $\mathcal{F}^\bullet$ of flat $\mathcal{O}$-modules.
Then $\mathcal{F}_0^\bullet =
\mathcal{F}^\bullet \otimes_{\mathcal{O}} \mathcal{O}_0$
represents $K_0$ in $D(\mathcal{O}_0)$.
Since $d_{\mathcal{F}^\bullet} \circ d_{\mathcal{F}^\bullet} = 0$
as $\mathcal{F}^\bullet$ is a complex, we see from the very construction
of $\omega(K_0)$ that it is zero.

\medskip\noindent
Assume (1). Let $\mathcal{F}^n$, $d : \mathcal{F}^n \to \mathcal{F}^{n + 1}$
be as in the construction of $\omega(K_0)$. The nullity of
$\omega(K_0)$ implies that the map
$$
\omega = d \circ d : \mathcal{F}_0^\bullet
\longrightarrow
(\mathcal{F}_0^\bullet \otimes_{\mathcal{O}_0} \mathcal{I})[2]
$$
is zero in $D(\mathcal{O}_0)$. By definition of the derived category
as the localization of the homotopy category of complexes
of $\mathcal{O}_0$-modules, there exists a quasi-isomorphism
$\alpha : \mathcal{G}_0^\bullet \to \mathcal{F}_0^\bullet$
such that there exist $\mathcal{O}_0$-modules maps
$h^n : \mathcal{G}_0^n \to
\mathcal{F}_0^{n + 1} \otimes_\mathcal{O} \mathcal{I}$
with
$$
\omega \circ \alpha =
d_{\mathcal{F}_0^\bullet \otimes \mathcal{I}} \circ h +
h \circ d_{\mathcal{G}_0^\bullet}
$$
We set
$$
\mathcal{H}^n = \mathcal{F}^n \times_{\mathcal{F}^n_0} \mathcal{G}_0^n
$$
and we define
$$
d' : \mathcal{H}^n \longrightarrow \mathcal{H}^{n + 1},\quad
(f^n, g_0^n) \longmapsto (d(f^n) - h^n(g_0^n), d(g_0^n))
$$
with obvious notation using that
$\mathcal{F}_0^{n + 1} \otimes_{\mathcal{O}_0} \mathcal{I} =
\mathcal{F}^{n + 1} \otimes_\mathcal{O} \mathcal{I} =
\mathcal{I}\mathcal{F}^{n + 1} \subset \mathcal{F}^{n + 1}$.
Then one checks $d' \circ d' = 0$ by our choice of $h^n$
and definition of $\omega$.
Hence $\mathcal{H}^\bullet$ defines an object in $D(\mathcal{O})$.
On the other hand, there is a short exact sequence of complexes
of $\mathcal{O}$-modules
$$
0 \to \mathcal{F}_0^\bullet \otimes_{\mathcal{O}_0} \mathcal{I} \to
\mathcal{H}^\bullet \to \mathcal{G}_0^\bullet \to 0
$$
We still have to show that
$\mathcal{H}^\bullet \otimes_\mathcal{O}^\mathbf{L} \mathcal{O}_0$
is isomorphic to $K_0$.
Choose a quasi-isomorphism
$\mathcal{E}^\bullet \to \mathcal{H}^\bullet$
where $\mathcal{E}^\bullet$ is a bounded above complex of flat
$\mathcal{O}$-modules. We obtain a commutative diagram
$$
\xymatrix{
0 \ar[r] &
\mathcal{E}^\bullet \otimes_\mathcal{O} \mathcal{I} \ar[d]^\beta \ar[r] &
\mathcal{E}^\bullet \ar[d]^\gamma \ar[r] &
\mathcal{E}_0^\bullet \ar[d]^\delta \ar[r] &
0 \\
0 \ar[r] &
\mathcal{F}_0^\bullet \otimes_{\mathcal{O}_0} \mathcal{I} \ar[r] &
\mathcal{H}^\bullet \ar[r] &
\mathcal{G}_0^\bullet \ar[r] &
0
}
$$
We claim that $\delta$ is a quasi-isomorphism. Since $H^i(\delta)$
is an isomorphism for $i \gg 0$, we can use descending induction
on $n$ such that $H^i(\delta)$ is an isomorphism for $i \geq n$.
Observe that
$\mathcal{E}^\bullet \otimes_\mathcal{O} \mathcal{I}$
represents
$\mathcal{E}_0^\bullet \otimes_{\mathcal{O}_0}^\mathbf{L} \mathcal{I}$,
that
$\mathcal{F}_0^\bullet \otimes_{\mathcal{O}_0} \mathcal{I}$
represents
$\mathcal{G}_0^\bullet \otimes_{\mathcal{O}_0}^\mathbf{L} \mathcal{I}$,
and that
$\beta = \delta \otimes_{\mathcal{O}_0}^\mathbf{L} \text{id}_\mathcal{I}$
as maps in $D(\mathcal{O}_0)$. This is true because
$\beta =
(\alpha \otimes \text{id}_\mathcal{I})
\circ
(\delta \otimes \text{id}_\mathcal{I})$.
Suppose that $H^i(\delta)$ is an isomorphism in degrees $\geq n$.
Then the same is true for $\beta$ by what we just said
and Lemma \ref{lemma-induced-map}.
Then we can look at the diagram
$$
\xymatrix{
H^{n - 1}(\mathcal{E}^\bullet \otimes_\mathcal{O} \mathcal{I})
\ar[r] \ar[d]^{H^{n - 1}(\beta)} &
H^{n - 1}(\mathcal{E}^\bullet) \ar[r] \ar[d] &
H^{n - 1}(\mathcal{E}_0^\bullet) \ar[r] \ar[d]^{H^{n - 1}(\delta)} &
H^n(\mathcal{E}^\bullet \otimes_\mathcal{O} \mathcal{I})
\ar[r] \ar[d]^{H^n(\beta)} &
H^n(\mathcal{E}^\bullet) \ar[d] \\
H^{n - 1}(\mathcal{F}_0^\bullet \otimes_\mathcal{O} \mathcal{I}) \ar[r] &
H^{n - 1}(\mathcal{H}^\bullet) \ar[r] &
H^{n - 1}(\mathcal{G}_0^\bullet) \ar[r] &
H^n(\mathcal{F}_0^\bullet \otimes_\mathcal{O} \mathcal{I}) \ar[r] &
H^n(\mathcal{H}^\bullet)
}
$$
Using Homology, Lemma \ref{homology-lemma-four-lemma}
we see that $H^{n - 1}(\delta)$ is surjective.
This in turn implies that $H^{n - 1}(\beta)$ is surjective
by Lemma \ref{lemma-induced-map}.
Using Homology, Lemma \ref{homology-lemma-four-lemma}
again we see that $H^{n - 1}(\delta)$ is an isomorphism.
The claim holds by induction, so $\delta$ is a quasi-isomorphism
which is what we wanted to show.
\end{proof}

\bibliographystyle{plain}
\bibliography{my}
\end{document}
